\subsubsection{Infraestructura - El entorno para ejecutar nuestras simulaciones [Pendientes: 1. Confirmar acceso, cuotas, almacenamiento y tarifas HPC-ZINE. 2. Verificar versión, instalación, configuración y ejecución efectiva de FDS en HPC-Zine. 3. Elegir sistema operativo y entornos de desarrollo de la aplicación. 4. Investigar y precisar los procedimientos, requisitos y actividades necesarios para obtener acceso individual a HPC-Zine.]}\label{sec:spmp-9-2-2}
Para desarrollar y ejecutar las simulaciones del proyecto, se prevé utilizar HPC-Zine como infraestructura universitaria de cómputo científico.

Juan Guillermo Torres Hurtado es la persona de contacto para la coordinación de HPC-ZINE. La infraestructura científica comprende entornos de ejecución de FDS, cómputo y almacenamiento.

El equipo ha completado la capacitación de HPC-Zine. Para obtener el acceso oficial a la infraestructura queda pendiente elaborar la cotización correspondiente. Cada integrante será responsable de su acceso independiente y de su utilización.
